% ========================================================================================
% INTRODUZIONE
% ========================================================================================
\chapter*{Introduzione}
\addcontentsline{toc}{chapter}{Introduzione}
\markboth{Introduzione}{}

Questo elaborato raccoglie la progettazione, l'implementazione in VHDL, la simulazione e,
dove richiesto, la sintesi su scheda di sviluppo degli esercizi proposti nel corso di
\thecorso. \todo{Breve descrizione del lavoro svolto e di eventuali parti opzionali (9 CFU).}

\section*{Strumenti e ambiente di sviluppo}

\begin{itemize}[nosep]
	\item \textbf{Linguaggio:} VHDL-93, con le librerie \sig{ieee.std\_logic\_1164} e \sig{ieee.numeric\_std}.
	\item \textbf{Ambiente:} AMD Vivado ML Standard 2023.1 (sintesi, implementazione, simulatore XSim).
	\item \textbf{Scheda:} Digilent Nexys A7-100T con FPGA Artix-7
	      (\texttt{xc7a100tcsg324-1}), clock di sistema a \SI{100}{\mega\hertz}.
\end{itemize}

\section*{Struttura di ogni esercizio}

Ogni esercizio segue la stessa struttura:
\begin{description}[nosep, font=\sffamily\bfseries]
	\item[Progetto e architettura] approccio di progetto, schema a blocchi del sistema e dei componenti, funzionalità.
	\item[Implementazione] codice VHDL dei componenti significativi; i componenti elementari riusati sono in appendice.
	\item[Simulazione] testbench principali e discussione dei risultati.
	\item[Sintesi su board di sviluppo] architettura aggiuntiva per l'I/O e file dei vincoli (se richiesta).
	\item[Timing analysis] analisi dei percorsi critici (se richiesta).
\end{description}

\section*{Convenzioni grafiche}

Nel documento si usano i seguenti riquadri.

\begin{traccia}
	Il testo ufficiale dell'esercizio, riportato senza modifiche.
\end{traccia}

\begin{scelta*}[Titolo]
	Una decisione di progetto lasciata aperta dalla traccia, con la sua motivazione.
	Nei capitoli le scelte sono numerate e richiamabili nel testo
	(es.\ ``come stabilito nella scelta~\ref{sc:mux-dest}'').
\end{scelta*}

\begin{osservazione}
	Un commento, un confronto o un dettaglio che aiuta a capire il progetto.
\end{osservazione}

\begin{attenzione}
	Un errore tipico, un caso limite o un vincolo da rispettare.
\end{attenzione}

\begin{risultato}
	L'esito di una simulazione o di una prova su scheda.
\end{risultato}

Nel testo, i nomi di segnali e porte sono scritti come \sig{clk} o \sig{d(3 downto 0)},
i nomi delle entity come \ent{mux16}.

\section*{Organizzazione dei sorgenti}

I sorgenti VHDL e i file di vincoli riportati nel documento sono inclusi direttamente dai
file usati per la simulazione e la sintesi, così il codice stampato coincide sempre con quello verificato.

\begin{file}[elaborato/]
\begin{verbatim}
main.tex, structure.tex
capitoli/          un file .tex per capitolo
codice/
    comuni/        componenti riusati (appendice)
    es1/ ... es12/ sorgenti, testbench e vincoli
figure/
    screenshot/    forme d'onda, report, foto
\end{verbatim}
\end{file}

Per compilare il documento:

\begin{commandline}
\begin{verbatim}
$ latexmk -pdf main.tex
$ ..\clean.ps1          # rimuove i file temporanei di LaTeX
\end{verbatim}
\end{commandline}
